\section{Implementation, verification, and computational comparisons}\label{sec:computation}
% Generated by verification/stage06-tables.py; data SHA256 373ce70b3839c12a206b193e9952d856f585da400f902bda773a725a609f73ac
\newcommand{\BenchPython}{3.12.14}
\newcommand{\BenchNumpy}{2.5.2}
\newcommand{\BenchScipy}{1.18.1}
\newcommand{\BenchHullObjective}{-6.981917}
\newcommand{\BenchMcObjective}{-7.287862}
\newcommand{\BenchBudgetObjective}{-6.538573}
\newcommand{\BenchCutCount}{42}
\newcommand{\BenchOldCold}{31.05}
\newcommand{\BenchNewThreeCold}{10.78}
\newcommand{\BenchNewThreeBases}{9.15}


The implementations serve two distinct purposes: constructing small numerical
LPs and producing exact rational hull certificates. The experiments compare
both with strengthened versions of classical disaggregation. They show large
model-size savings, but also cases where a simpler LP is faster and cases where
global state merging explains most of an apparent compression advantage.

\subsection{Implemented scope and certificate semantics}

\begin{table}[tb]\centering\small
\caption{Implemented scope. Exact assembly of an LP does not make its
floating-point optimizer an exact certificate.}\label{tab:implemented-scope}
\begin{tabular}{>{\raggedright\arraybackslash}p{.19\linewidth}>{\raggedright\arraybackslash}p{.36\linewidth}>{\raggedright\arraybackslash}p{.36\linewidth}}\toprule
Component & Supported input & Returned evidence\\\midrule
Parallel-path separator & Arbitrarily oriented multigraphs whose cyclic
blocks are parallel paths; loops and bridges included & Exact original-coordinate
cut or compact rational decomposition\\[3pt]
Flat-chain separator & Canonical unit-data $G_L$, arbitrary observations;
parameter is the number of observed labels & Exact cut or compact decomposition;
unit flow/product coefficients through three observed labels\\[3pt]
General compressed EF & Arbitrary bounded equality-flow graphs; initial
or observed-coordinate-eliminated formulation & Exact rational model;
numerical LP solution; optionally an exactly checked Farkas cut\\[3pt]
Bounded-rank libraries & Supplied small signed-normal systems and structural
verification instances & Verification prototype for support and recovery;
not a complete production general-graph oracle\\\bottomrule
\end{tabular}
\end{table}

\Cref{tab:implemented-scope} distinguishes the three usable interfaces
from the bounded-rank verification code. The specialized separators use
Python's standard library and rational arithmetic throughout. They check
original flow and simplex constraints, retain the original affine expressions
of active bounds, and return a cut with strictly positive exact violation
or a decomposition whose positive-weight normalized flows belong to $\Flow$.
The graph separator rejects unsupported blocks explicitly. The flat oracle
merges unobserved labels, eliminates the residual profile, and uses the
five-test formula when at most two labels are observed. At three labels it
uses the 16 circuits and both flow-balance repairs of
\cref{thm:three-state-chain} for separation. Recovery uses finite inverse
bases at three or more observed labels.
Its stored default is shared across unused global labels. Compatibility
accessors can materialize the former dense state arrays, with the corresponding
output cost. Tuple normal construction and hashing retain a dimension factor
relative to the preindexed-mask operation count in \cref{thm:fixed-chain};
apart from the constructor's $O(|\Obs|\log|\Obs|)$ observation sorting,
the implemented cost is linear in $L$ at fixed observed-label count.

The general formulation stores exact rational coefficients, including the
observed-coordinate substitutions, but calls a numerical LP solver.
Its phase-I certificate routine rationalizes dual weights and, if necessary,
repairs their supported kernel equations by exact elimination. A returned
cut is accepted only after checking nonnegative weights, exact cancellation
of every auxiliary coordinate, and strictly positive rational violation.
The API distinguishes \texttt{certified\_outside},
\texttt{numerically\_feasible}, \texttt{uncertified\_outside}, and
\texttt{solver\_failure}. The second status is not exact membership,
and failure of exact multiplier recovery is not an infeasibility certificate.
This is an implementation of classical Farkas projection, not a new generic
separation theorem.

Code indices are zero based, with explicit states $0,\ldots,m-1$ and
residual state $m$; this is only an indexing change from the manuscript.
The LP baseline generator uses outgoing-minus-incoming balances, and its
adapter negates them before calling either production model. Float inputs
to exact interfaces mean their decimal text. The optimization experiment
uses a small-denominator reconstruction adapter, requires error below
$10^{-7}$ and exact aggregate balances, and then checks the returned
decomposition exactly. This adapter is suitable for the stated synthetic
instances; it is not a general certified rational LP optimizer.

The verification suite compares against independently assembled full-state
LPs and, for small graphs, directly enumerated path--simplex vertices.
The flat tests include arbitrary observations and zero weights,
both three-label repairs with exact global cut checks, the four-label
ratio-two obstruction, and sparse observed labels among many original
simplex coordinates. They also test an interior-residual point with zero
explicit profile and positive residual branch flow: this rules out
incorrectly reusing a basis construction that forces the explicit profile
sum to the full branch flow. Separate tests check every inverse basis for recovery with three
observed labels. An independent flat audit verifies 192 exact
decompositions and 708 exact violated cuts, with 233 numerical path-hull
comparisons. The general compressed audit passes 600 objective comparisons,
758 membership comparisons, and 476 exactly checked Farkas cuts, each also
tested by numerical full-hull support optimization. Its structural integration
checks recover the coefficient ratios two, five, and twenty-one. These finite
checks support the code; the preceding proofs establish the general claims.

\subsection{Baselines that retain the elementary reductions}

For optimization with fixed weights, the full disaggregation uses only
positive-weight state flows, block balances $Af^j=\lambda_j b$, and native
variable bounds $0\le f^j\le\lambda_j u$. It eliminates $x,z$ by substitution
and adds the constant $c_y^\top\bar y$ to the objective value. In additional
original-coordinate rows, the fixed $y$ contribution moves to the right-hand
side. This is one joint LP and remains valid with the aggregate budget below.
The global-merger baseline first groups labels observed nowhere into one
default, then applies the same construction to its positive-weight states.
With free weights, both interfaces retain every original $y$ variable and its
objective and constraint contributions. Both use sparse block and Kronecker
assembly. They use no graph blocks, cycle coordinates, or observation
rank elimination. Initial compression and observed elimination use the two
variants of the general implementation, including their preprocessing costs.

For point membership, $x,y,z$ are data and only state flows are LP variables.
The positive-state full baseline removes exactly zero-weight states, checks their observations, and
assembles state balances and aggregate equations as sparse block matrices.
It accepts solver statuses zero and two as numerical feasibility and
infeasibility; other statuses are failures, not successful rejections.

An additional elementary baseline applies when precisely two states have
positive weights $\lambda,\mu$, with $\lambda+\mu=1$. After checking
the original aggregate flow, write the state flows as $f$ and $x-f$.
The complete remaining feasibility system is
\begin{equation}\label{eq:two-state-membership-lp}
 Af=\lambda b,\qquad
 \max\{0,x-\mu u\}\le f\le\min\{\lambda u,x\}.
\end{equation}
An observation in the first state fixes $f_e=z_{ej}$; an observation
in the second fixes $f_e=x_e-z_{ej}$. Intersect these fixed values
with both bounds and with each other. The second state's balance follows
from $Ax=b$, so this one-flow network LP is exact as a formulation.
One positive state requires only a direct check. In the numerical implementation
the aggregate balance precheck uses tolerance $10^{-9}$; its LP answer is
still numerical evidence. Exact input arithmetic identifies zero states and
checks the intersected bounds.

Fixed-weight linear optimization without additional coupling is also simpler
than a joint LP suggests. If $c^j_e$ is the product-objective coefficient
for $(e,j)$, or zero when absent, and $c^0=0$, then disaggregation gives
\begin{equation}\label{eq:independent-network-objectives}
 \min_{(x,\bar y,z)\in\Hull}
       (c_x^\top x+c_y^\top\bar y+c_z^\top z)
 =c_y^\top\bar y+
   \sum_{j\in\states}\lambda_j
       \min_{v\in\Flow}(c_x+c^j)^\top v.
\end{equation}
This follows by independently minimizing each positive-weight state flow;
zero-weight terms vanish. Unobserved labels have identical costs and share
one solve. We include this sequential network-LP baseline. With free weights
and no side constraints, a simplex vertex suffices, so such stand-alone
objectives should not be presented as intrinsically difficult coupled problems.

\subsection{Instances and measurement protocol}

All reported data are synthetic. The code uses SciPy
\citep{Virtanen2020} and its bundled HiGHS LP solver
\citep{HuangfuHall2018}, with default \texttt{highs} options.
The recorded versions are Python~\BenchPython, NumPy~\BenchNumpy,
and SciPy~\BenchScipy, on x86-64 Linux under WSL2 with 36 visible
logical CPUs. The bundled HiGHS version is 1.12.0. No affinity, host
isolation, persistent model, or solver warm-start advantage is claimed.
Each case has one recorded warmup followed by five runs with method order
rotated. Each run rebuilds its model. Build and solver times are recorded
separately, together with total times and all minimum/median/maximum values.
Component medians need not add to the total median. Independent certificate
audits occur outside the timed components and have their own recorded cost.
Membership-only and decomposition-inclusive exact queries are separate
workloads, not timings inferred by subtraction.

The principal sparse optimization instance has four articulation-linked
blocks, each containing five internally disjoint paths of length two,
with a bridge between successive blocks. It has 43 arcs and 128 explicit
labels. Capacities are integers from two to eight, orientations vary, and
balances are induced by the half-capacity reference flow. Each block observes
three labels on three arcs each. The graph and label-generation seeds are
fixed in the accompanying generator; the primary seed is 33, and objective
seed 1033 gives independent normal costs with standard deviation $0.15$
on $x$, zero on $y$, and standard deviation one on observed $z$.
Weights are fixed to $9/(10m)$.

The membership family uses four six-path blocks of length two, seed 41,
51 arcs, 36 observations, and $m=16,128,1024$. Each explicit state has
weight $1/(2m)$, leaving residual weight $1/2$. In block $B$, state
$j$ perturbs path $(j+B)\bmod6$ and the next path by opposite amounts
$(j\bmod3-1)/4$, with path orientation signs applied. The aggregate
and observations are formed from this exact mixture around the half-capacity
reference. These are feasible membership queries, not random rejection tests.

Two controls separate the mechanisms. First, an all-labels-observed instance
has 16 three-path blocks, four disjoint labels per block, and one observation
on the first arc of each path for each local label. All 64 labels occur
somewhere, so global merging cannot help. It has 111 arcs and 192 observations;
graph seed 73 and objective seed 1073 are fixed, with nonzero $y$ costs and
weights $9/640$. Second, on the original 43-arc instance, choose the
nonnegative resource indicator $d_e=1$ where an unconstrained optimum has
$x_e$ above the reference value, and zero elsewhere. Put
\[
 d^\top x\le B,\qquad
 B=\tfrac12(d^\top x^{\rm opt}+d^\top x^{\rm ref}).
\]
The exact reconstructed unconstrained point violates this row, while the
reference graph point satisfies it. Every formulation receives the same
aggregate row. This compares equivalent relaxations $\Hull\cap\{d^\top x\le B\}$;
it does not assert that convexification commutes with adding a budget.
The independent-state baseline \eqref{eq:independent-network-objectives}
is inapplicable to this coupled control.

Finally the flat family uses $G_L$ for $L=8,32,128,512$, observing
$(a_i,1)$ and $(a_i,2)$ alternately. Boundary weights are $(1/2,1/2)$;
interior weights are $(1/3,1/3)$, so the latter has three positive
states. Feasible queries have $x_{a_i}=x_{b_i}=1/4$, $x_h=1/2$.
Writing $g=i-1$, the boundary observed value is
$1/8+(-1)^g(g\bmod3-1)/16$; the interior value is
$1/12+(g\bmod3-1)/48$. Branch profiles $1/4,1/4$ in the boundary
case and $1/6,1/6,1/6$ in the interior case give feasible completions:
in the latter, split the remaining $a_i$ aggregate equally between
the other two states. Infeasible queries use
$x_{a_i}=3/10$, $x_{b_i}=1/5$, $x_h=1/2$, and every observed
product $3/10$. All separate McCormick inequalities hold, but the
two observed profiles would sum to at least $3/5>1/2$.

\subsection{Size savings and runtime tradeoffs}

% Generated by verification/stage06-tables.py; data SHA256 373ce70b3839c12a206b193e9952d856f585da400f902bda773a725a609f73ac
\begin{table}[tb]\centering\footnotesize
\caption{The 43-arc, 128-label optimization case. Times are milliseconds; total reports median [minimum,maximum] over five runs. The network baseline reuses one matrix for sequential solves. Repeated-cut build/solve columns exclude its separately recorded separator and reconstruction work, which the total includes.}
\label{tab:computation-opt}
\setlength{\tabcolsep}{3pt}
\begin{tabular}{lrrr rrr}\toprule
Method & Vars. & Rows & Nonzeros & Total & Build & Solve\\\midrule
Full & 5547 & 3612 & 11094 & 19.05 [17.81,23.05] & 1.07 & 18.12 \\
Global merge & 516 & 336 & 1032 & 4.25 [3.41,5.01] & 0.70 & 3.18 \\
Initial compression & 255 & 225 & 936 & 4.82 [4.70,5.62] & 2.48 & 2.31 \\
Observed elimination & 222 & 192 & 839 & 6.45 [5.52,9.31] & 4.11 & 2.47 \\
Independent networks & 43 & 28 & 86 & 22.02 [19.39,28.09] & 0.37 & 21.64 \\
Repeated cuts & 207 & 179 & 630 & 292.22 [280.98,330.97] & 14.78 & 100.57 \\
\bottomrule\end{tabular}\end{table}

% Generated by verification/stage06-tables.py; data SHA256 373ce70b3839c12a206b193e9952d856f585da400f902bda773a725a609f73ac
\begin{table}[tb]\centering\small
\caption{Feasible membership on the 51-arc, 36-observation family: median construction-plus-query milliseconds. The exact-separator column is membership without requested decomposition; the raw data separately record constructive queries and all timing ranges.}
\label{tab:computation-member}
\begin{tabular}{rrrrrr}\toprule
Labels & Full & Global merge & Initial & Eliminated & Exact separator\\\midrule
16 & 7.42 & 5.35 & 6.21 & 7.65 & 5.93 \\
128 & 44.65 & 5.01 & 5.54 & 7.14 & 5.23 \\
1024 & 1580.67 & 9.29 & 16.16 & 16.25 & 13.68 \\
\bottomrule\end{tabular}\end{table}

% Generated by verification/stage06-tables.py; data SHA256 373ce70b3839c12a206b193e9952d856f585da400f902bda773a725a609f73ac
\begin{table}[tb]\centering\small
\caption{Two controls: median total optimization milliseconds. Independent network optimization is inapplicable with the aggregate budget. Repeated cutting was measured on the original and budgeted instance, not on the all-labels-observed control.}
\label{tab:computation-ablations}
\setlength{\tabcolsep}{3pt}
\begin{tabular}{lrrrrrr}\toprule
Case & Full & Global & Initial & Eliminated & Networks & Cuts\\\midrule
Aggregate budget & 32.80 & 4.33 & 5.84 & 6.11 & -- & 273.35 \\
All labels observed & 30.36 & 30.38 & 10.16 & 18.39 & 129.00 & -- \\
\bottomrule\end{tabular}\end{table}


\Cref{tab:computation-opt} shows that global merging accounts for much
of the full-versus-compressed difference on the original sparse instance.
Only 11 of its 128 labels are observed anywhere: the full formulation's
5,547 variables drop to 516 under global merging, then to 255 under initial
compression and 222 under observed elimination. Elimination removes 33 of
48 auxiliary coordinates. It saves rows and nonzeros but is slower overall
than initial compression in this run: its extra symbolic work is not offset
by a solver saving. The globally merged and initially compressed
LPs have median total times of 4.25 and 4.82 milliseconds, respectively;
global merging has the lower median despite its larger model, although the
five-run timing ranges overlap.
Comparing only with the full formulation would overstate the specifically
graph-local benefit.

The McCormick optimum is $\BenchMcObjective$, compared with hull optimum
$\BenchHullObjective$. The repeated single-cut loop requires
\BenchCutCount\ cuts and is slower than the compact LPs. Its final
decomposition is checked exactly, and each warmup cut receives an independent
numerical global support audit. The independent-network baseline also attains
the same objective, but multiple solver calls have overhead; algebraic
separability does not guarantee the fastest implementation.

\Cref{tab:computation-member} gives the analogous membership comparison.
At 1,024 labels, global merging is faster than either general compressed
variant and the exact separator on these instances. The full LP's much
larger positive-state system is unnecessary when almost all labels are
unobserved. The exact separator supplies a rational constructive certificate
when requested, which is a different output from numerical LP feasibility.

The all-labels-observed control in \cref{tab:computation-ablations}
isolates a benefit that global merging cannot provide. Both full and global
formulations have 7,215 variables; initial compression has 495, and observed
elimination has 367 with no auxiliary variables. Initial compression is
substantially faster here. The aggregate-budget control retains the same
size ordering and has objective $\BenchBudgetObjective$. Global merging is
also faster than initial compression on that control, with median totals of
4.33 and 5.84 milliseconds. Both remain faster than repeated cutting on this
coupled relaxation.
The budget and the local-label control are small targeted comparisons, not
evidence about industrial instances.

% Generated by verification/stage06-tables.py; data SHA256 373ce70b3839c12a206b193e9952d856f585da400f902bda773a725a609f73ac
\begin{table}[tb]\centering\small
\caption{Flat-chain membership: median construction-plus-query milliseconds. Boundary weights are $(1/2,1/2)$; interior weights are $(1/3,1/3)$. Full retains only positive states. The one-flow reduction applies only with two positive states. Exact columns exclude optional decomposition; all ranges and constructive workloads are retained in the data.}
\label{tab:computation-flat}
\begin{tabular}{rllrrrr}\toprule
$L$ & Weights & Point & Full & One-flow LP & Unreduced exact & Reduced exact\\\midrule
8 & Boundary & In & 2.09 & 1.60 & 0.44 & 0.33 \\
8 & Boundary & Out & 1.89 & 1.36 & 0.49 & 0.37 \\
8 & Interior & In & 2.53 & -- & 0.46 & 0.31 \\
8 & Interior & Out & 2.90 & -- & 0.79 & 0.51 \\
32 & Boundary & In & 3.09 & 2.12 & 1.38 & 0.97 \\
32 & Boundary & Out & 3.65 & 1.83 & 1.53 & 1.07 \\
32 & Interior & In & 3.46 & -- & 1.45 & 1.00 \\
32 & Interior & Out & 2.48 & -- & 1.32 & 1.04 \\
128 & Boundary & In & 4.95 & 4.41 & 5.06 & 3.85 \\
128 & Boundary & Out & 5.63 & 4.57 & 6.23 & 4.03 \\
128 & Interior & In & 10.39 & -- & 4.53 & 3.92 \\
128 & Interior & Out & 4.12 & -- & 5.34 & 4.11 \\
512 & Boundary & In & 14.36 & 12.90 & 18.89 & 16.93 \\
512 & Boundary & Out & 10.62 & 9.31 & 20.37 & 15.74 \\
512 & Interior & In & 41.04 & -- & 27.16 & 23.12 \\
512 & Interior & Out & 39.18 & -- & 51.05 & 41.23 \\
\bottomrule\end{tabular}\end{table}


The elementary reductions matter on long chains
(\cref{tab:computation-flat}). On the 512-gadget boundary face, both
the positive-state block LP and the one-flow LP are faster than the reduced
exact oracle. The positive-state LP treats original coordinates as data,
removes zero-weight states, and uses sparse block assembly. The one-flow
substitution gives a further reduction. At that boundary point,
the positive-state LP has 2,050 variables and the one-flow LP has 1,025;
the interior three-state LP has 3,075.
On feasible interior queries the reduced oracle is faster at the largest
length. Infeasible interior queries have similar medians and wide ranges:
the full LP ranges from 16.08 to 45.27 milliseconds and the reduced oracle
from 21.02 to 58.95 milliseconds. Five measurements on a shared host do
not support a universal speed ordering.

The reduced oracle merges unobserved labels and eliminates the default
profile coordinate before separation. The unreduced comparator retains
that coordinate and uses a signed-subset circuit library; it supplies the
same exact membership and decomposition outputs. The two-observed-label
reduced oracle uses five explicit tests and has no circuit-library setup.
For comparison, the unreduced three-coordinate library takes
\BenchOldCold\ milliseconds to construct in the recorded cold call.
The reduced three-observed-label library
takes \BenchNewThreeCold\ milliseconds, with another
\BenchNewThreeBases\ milliseconds for its inverse bases; these are separate
from the two-label workload. Initial and observed-eliminated general models
were measured on the eight flat cases at lengths eight and 32. They do not
improve the total times over the dedicated exact oracle on those cases;
their complete component and size data are included in the supplement.
The general compressed formulations are not timed at the two larger lengths
in this comparison.

Rows in the tables exclude variable bounds. Nonzero and byte counts refer
to input sparse matrices, not peak memory, solver factors, or Python rational
objects. These measurements establish the behavior of the supplied code on
the stated inputs. They neither compare optimized native implementations
nor measure persistent callbacks in a larger optimizer.

\subsection{A two-supplier service-transportation interpretation}

The service-conflict transportation model of
\citet[Section~4.2]{KhademniaDavarnia2025} motivates the following
two-supplier specialization with a common simplex.
Consider a transportation network with two suppliers
$s_1,s_2$ and demand nodes $v_1,\ldots,v_k$, with arcs from both suppliers
to every demand node. Let supplies be $S_1,S_2$, demands $D_i$, and
arc capacities $u_{\ell i}$. The flow equations are
\[
 \sum_i x_{\ell i}=S_\ell\quad(\ell=1,2),\qquad
 x_{1i}+x_{2i}=D_i\quad(i\in[k]),\qquad
 0\le x_{\ell i}\le u_{\ell i}.
\]
Suppose a service or contract allocation $y\in\Delta_m$ modifies the
cost on selected route--service pairs. Its objective contribution has the
form $\sum_{((\ell,i),j)\in\Obs}c_{\ell i,j}x_{\ell i}y_j$.
Only pairs with a nonzero adjustment need product coordinates.

The underlying graph is a union of $k$ internally disjoint length-two
paths between the suppliers, even though the second arc on each path is
traversed against its direction. It is therefore covered by
\cref{thm:parallel-hull}. Demand conservation means that a circulation
increase on the first supplier's arc is the opposite change on the second
supplier's arc. The path coordinates retain these identities, while the
state transportation problem enforces compatibility of all selected service
products. With three demand nodes the theta formulas apply directly.
For larger $k$, an exact transportation cut or compact EF is available;
the number of observed service labels, their path coverage, and any repeated
observations determine the useful reduction.

This interpretation uses the common balance and capacity domain assumed
throughout the paper. An outer resource budget or linking constraint can be
added to the exact component hull to strengthen a relaxation, as in the
budget experiment. The resulting intersection is not generally the convex
hull of the additionally constrained product graph. This is a modeling
illustration; the experiments above use synthetic networks.

\subsection{Reproduction}

The accompanying computational supplement contains the exact interfaces,
numerical formulations, independent checks, measurement generator, and
canonical raw data used in the tables. Its top-level README gives dependency
versions, certificate semantics, and commands for the regression suite and
independent checks. Paths in the commands below are relative to the extracted
supplement root. The preserved directory layout is required by the benchmark
runner's import of the unreduced flat-chain comparator.

To repeat the complete measurement protocol, write new results to a separate
file:
\begin{verbatim}
PYTHONPATH=code python -m network_simplex_benchmarks.paper_stage06 \
  --output rerun-benchmarks.json
\end{verbatim}
The default is one recorded warmup followed by five rotated measurements
per case. The reported optimization and membership measurements were collected
separately; all methods within an optimization case use identical objectives,
weights, and additional rows. Rerunning the generator repeats these inputs
and protocols, but timing values depend on the machine and software.
The supplied data preserve the reported measurements rather than replacing
them with new timings.

To regenerate the published tables from those canonical data, run
\begin{verbatim}
python paper-network-simplex/verification/stage06-tables.py
\end{verbatim}
The generator recomputes every reported minimum, median, and maximum from
the raw five-run records, verifies the complete case grid and method order,
and writes the five table and value files under
\texttt{paper-network-simplex/tables/}. The supplement manifest hashes the
runnable source, independent checks, raw measurements, and generated tables.
The data include input definitions, objective vectors, additional rows,
warmups, timing components, exact witness and cut audits, and numerical solver
checks. The standalone LaTeX source archive includes all manuscript inputs
and a separate build README.
